% 数学ノート集・共通テンプレート v2（積分論2026年度第1部の版面を参考）
% A5判。章内で定義・定理・例・注意等の番号を共通化。
% コンパイル： platex -interaction=nonstopmode math_note_template_v2.tex (2回)
%              dvipdfmx math_note_template_v2.dvi
\documentclass[a5paper,10pt,onecolumn]{jsarticle}
\usepackage[dvipdfmx,top=22mm,bottom=20mm,inner=19mm,outer=18mm,headheight=13pt,headsep=6mm]{geometry}
\usepackage{amsmath,amssymb,amsthm,mathtools}
\usepackage{enumitem,fancyhdr,titlesec}
\usepackage[dvipdfmx,bookmarks=true,colorlinks=false]{hyperref}
\usepackage{pxjahyper}

% ----- 差し替える文書情報 -----
\newcommand{\NoteTitle}{Haar測度と局所コンパクト群}
\newcommand{\NoteSeries}{数学ノート集}

% ----- 文字・組版 -----
% 参考PDFの質素な明朝＋Computer Modernの印象を引き継ぎながら，
% A5判で窮屈にならないよう本文の行間を少し確保する．
\linespread{1.13}
\setlength{\parindent}{1zw}
\setlength{\parskip}{0pt}
\setlength{\abovedisplayskip}{6pt plus 2pt minus 2pt}
\setlength{\belowdisplayskip}{6pt plus 2pt minus 2pt}
\setlength{\abovedisplayshortskip}{3pt plus 2pt minus 1pt}
\setlength{\belowdisplayshortskip}{4pt plus 2pt minus 1pt}
\setlist[enumerate]{topsep=3pt,itemsep=1pt,parsep=0pt,leftmargin=2.4zw}
\titleformat{\section}{\normalfont\normalsize\bfseries}{\thesection}{0.8zw}{}
\titlespacing*{\section}{0pt}{2.0ex plus .4ex}{.9ex}
\titleformat{\subsection}{\normalfont\normalsize\bfseries}{\thesubsection}{0.7zw}{}
\titlespacing*{\subsection}{0pt}{1.5ex plus .3ex}{.6ex}
\raggedbottom

% ----- 柱とページ番号：ページ上部、罫線なし -----
\fancypagestyle{mathnote}{%
  \fancyhf{}%
  \fancyhead[L]{\footnotesize \NoteSeries\ （\NoteTitle）}%
  \fancyhead[R]{\footnotesize \thepage}%
  \renewcommand{\headrulewidth}{0pt}%
  \renewcommand{\footrulewidth}{0pt}%
}
\pagestyle{mathnote}

% ----- 節内通し番号の定義・定理群 -----
\newtheoremstyle{mathnotetheorem}
 {4.5pt}{4.5pt}{\normalfont}{}%
 {\bfseries}{}{0.65zw}%
 {\thmname{#1}\thmnumber{ #2}\thmnote{（#3）}}
\theoremstyle{mathnotetheorem}
\newtheorem{theorem}{定理}[section]
\newtheorem{definition}[theorem]{定義}
\newtheorem{lemma}[theorem]{補題}
\newtheorem{proposition}[theorem]{命題}
\newtheorem{corollary}[theorem]{系}
\newtheorem{example}[theorem]{例}
\newtheorem{remark}[theorem]{注意}
\renewcommand{\proofname}{証明}
\renewcommand{\qedsymbol}{$\square$}
\makeatletter
\renewenvironment{proof}[1][\proofname]{%
  \par\pushQED{\qed}\normalfont\topsep4\p@\@plus2\p@\relax
  \trivlist\item[\hskip\labelsep\bfseries #1]\ignorespaces
}{\popQED\endtrivlist\@endpefalse}
\makeatother

% ----- 使うノートで適宜拡張する記号 -----
\newcommand{\R}{\mathbb{R}}
\newcommand{\norm}[1]{\lVert #1\rVert}
\newcommand{\abs}[1]{\lvert #1\rvert}

\begin{document}
\section*{読了後にできること}
本稿を読むための前提は，局所コンパクト群，Haar測度，左不変性．
本稿の具体的な到達課題は，具体的なコンパクト群でHaar測度を規格化して積分を計算する．


\par\noindent\textbf{記号と約束.}\quad
関数空間の添字 $c$ はコンパクトな台をもつことを表す．
\par

\section{局所コンパクト群とRadon測度}

本稿では $G$ を局所コンパクトHausdorff位相群とする．
測度論についてはRadon測度の外正則性，
開集合上の内正則性，Carath\'eodoryの外測度構成，
およびTonelliの定理を前提とする．
局所コンパクト空間のUrysohnの補題も用いる．
本稿では $C_c(G)$ を $G$ 上のコンパクトな台を持つ連続実数値関数の空間とする．
ここで関数の台は非零点集合の閉包である．

\begin{definition}
$G$ 上の正のBorel測度 $\mu$ が\emph{Radon測度}であるとは，
各コンパクト集合 $K$ で $\mu(K)<\infty$ となり，
任意のBorel集合 $A$ について
\[
 \mu(A)=\inf_{\substack{U\supset A\\U\ {\rm open}}}\mu(U)
\]
となり，任意の開集合 $U$ について
\[
 \mu(U)=\sup_{\substack{K\subset U\\K\ {\rm compact}}}\mu(K)
\]
となることをいう．
\end{definition}

\begin{definition}
零でないRadon測度 $\mu$ が
\[
 \mu(gA)=\mu(A)
 \qquad(g\in G,\ A\text{ はBorel集合})
\]
を満たすとき，\emph{左Haar測度}という．
右移動 $A\mapsto Ag$ に関して不変なRadon測度を
右Haar測度という．
\end{definition}

\begin{lemma}[コンパクト集合と平行移動]\label{lem:topological}
次が成り立つ．
\begin{enumerate}
\item $G$ の任意の非空開集合の左平行移動全体は $G$ を覆う．
\item コンパクト集合 $K$ と開集合 $U\supset K$ に対し，
 $K\subset V\subset\overline V\subset U$ となる開集合 $V$ で
 $\overline V$ がコンパクトなものを取れる．
\item 交わらない二つのコンパクト集合 $K,L$ に対し，
 単位元 $e$ の近傍 $W$ で，
 どの左平行移動 $gW$ も $K,L$ の両方とは交わらないものがある．
\end{enumerate}
\end{lemma}
\begin{proof}
(1) $U$ の元 $u$ を固定すると任意の $x\in G$ は
$(xu^{-1})u\in(xu^{-1})U$ に属する．
(2) 各 $x\in K$ で局所コンパクト正則性から
$x\in V_x\subset\overline{V_x}\subset U$，
$\overline{V_x}$ がコンパクトとなる近傍を取る．
有限個の $V_x$ が $K$ を覆うのでその和集合を $V$ とすればよい．
(3) $K^{-1}L$ はコンパクトで $e$ を含まない．
連続性から単位元の近傍 $W$ を十分小さく取れば
$W^{-1}W\cap K^{-1}L=\varnothing$ となる．
もし $gW$ が $K,L$ の双方に交われば
$k=gw_1\in K$，$\ell=gw_2\in L$ となり，
$k^{-1}\ell=w_1^{-1}w_2$ が上の空集合に属する．
矛盾するので主張が従う．
\end{proof}

\section{コンパクト集合の被覆数}

\begin{definition}
コンパクト集合 $K\subset G$ と
非空開集合 $U\subset G$ に対し，
\[
 N(K,U)=\min\left\{r\ge0:
 K\subset\bigcup_{j=1}^r g_jU
 \text{ となる }g_j\in G\text{ がある}\right\}
\]
と定義する．ただし $N(\varnothing,U)=0$ とする．
\end{definition}

\begin{lemma}[被覆数の基本性質]\label{lem:cover}
任意のコンパクト集合 $K,L$ と
非空開集合 $U$ について
\[
 N(K,U)<\infty,\qquad
 N(gK,U)=N(K,U),
\]
\[
 K\subset L\ \Longrightarrow\ N(K,U)\le N(L,U),
\]
\[
 N(K\cup L,U)\le N(K,U)+N(L,U)
\]
が成り立つ．
さらに $K\cap L=\varnothing$ なら，
十分小さい単位元の近傍 $U$ について
\[
 N(K\cup L,U)=N(K,U)+N(L,U)
\]
となる．
\end{lemma}
\begin{proof}
補題\ref{lem:topological}(1)と $K$ のコンパクト性から
有限個の左平行移動が $K$ を覆う．
左平行移動は被覆全体を同じ個数の平行移動へ写すので
被覆数を変えない．
単調性と劣加法性は被覆の定義から従う．
$K,L$ が交わらなければ
補題\ref{lem:topological}(3)の近傍 $W$ を選ぶ．
$U\subset W$ とし，$K\cup L$ の最小被覆を考える．
各 $g_jU$ は $K$ と $L$ の両方には交わらないので，
被覆のうち $K$ に交わるものだけで $K$ を，
$L$ に交わるものだけで $L$ を覆える．
従って被覆数は
$N(K,U)+N(L,U)$ 以上となる．
逆向きの不等式は劣加法性による．
\end{proof}

\begin{lemma}[正規化被覆数の有界性]\label{lem:bounded-ratios}
空でない内部を持つコンパクト集合 $K_0$ を固定する．
任意のコンパクト集合 $K$ に対して，
$K$ のみに依存する整数 $m_K$ が存在して
\[
 0\le\frac{N(K,U)}{N(K_0,U)}\le m_K
\]
が任意の単位元の開近傍 $U$ について成立する．
\end{lemma}
\begin{proof}
$K_0^\circ$ は非空開集合なので，
その左平行移動全体は $G$ を覆う．
$K$ のコンパクト性から
$K\subset\bigcup_{j=1}^{m_K}g_jK_0^\circ
\subset\bigcup_{j=1}^{m_K}g_jK_0$
となる有限個の $g_j$ を取れる．
補題\ref{lem:cover}の単調性・劣加法性・移動不変性から
\[
 N(K,U)\le\sum_{j=1}^{m_K}N(g_jK_0,U)
 =m_KN(K_0,U).
\]
$K_0$ は空でないので分母は1以上であり，結論が従う．
\end{proof}

\section{不変なコンパクト内容の構成}

\begin{theorem}[不変なコンパクト内容]\label{thm:content}
コンパクト集合全体の上に実数値写像 $c$ が存在して，
次を満たす．
\begin{enumerate}
\item $c(\varnothing)=0$，$c(K_0)=1$，
  $0\le c(K)<\infty$．
\item $K\subset L$ なら $c(K)\le c(L)$．
\item $c(gK)=c(K)$（$g\in G$）．
\item $c(K\cup L)\le c(K)+c(L)$．
\item $K\cap L=\varnothing$ なら
$c(K\cup L)=c(K)+c(L)$．
\end{enumerate}
\end{theorem}
\begin{proof}
単位元の開近傍全体を，包含の逆順によって有向集合とする．
各近傍 $U$ に対して
\[
 c_U(K)=\frac{N(K,U)}{N(K_0,U)}
\]
と定める．
補題\ref{lem:bounded-ratios}によって
$c_U(K)\in[0,m_K]$ である．
コンパクト集合全体を添字とする直積
$\prod_K[0,m_K]$ はTychonoffの定理によってコンパクトである．
従ってネット $(c_U)_U$ はこの直積空間に収束する
部分ネットを持つ．その座標ごとの極限を $c(K)$ と置く．
各 $c_U$ は非負，正規化され，
単調，移動不変，劣加法的であるので，
不等式の極限から (1)--(4) が成立する．
$K,L$ が交わらなければ，補題\ref{lem:cover}によって
十分小さい近傍 $U$ で
$c_U(K\cup L)=c_U(K)+c_U(L)$ が成立する．
収束部分ネットは近傍系を共終的に細かくするので，
その極限でも等式が成立する．従って (5) を得る．
\end{proof}

\section{コンパクト内容から測度へ}

\begin{lemma}[Radon測度への拡張]\label{lem:content-extension}
$X$ を局所コンパクトHausdorff空間とする．
コンパクト集合上の写像 $c$ が
定理\ref{thm:content}の非負性・有限性・単調性・
劣加法性・互いに素なコンパクト集合への加法性を
満たすとする．
開集合 $U$ に対し
\[
 m(U)=\sup_{K\subset U,\ K\ {\rm compact}}c(K)
\]
と置く．このとき
\[
 \mu^*(A)=\inf_{A\subset U,\ U\ {\rm open}}m(U)
\]
は外測度であり，そのBorel集合への制限は
Radon測度を与える．
さらに $\mu(U)=m(U)$ が全ての開集合で成立する．
\end{lemma}
\begin{proof}
$m$ は開集合の包含について単調であり，
$m(\varnothing)=0$ である．
開集合族 $(U_i)$ を取り，
コンパクト集合 $K\subset\bigcup_iU_i$ を固定する．
コンパクト性によって有限個の $U_i$ が $K$ を覆う．
局所コンパクトHausdorff空間の正則性を用いて，
コンパクト集合 $K_i\subset U_i$ で
$K\subset\bigcup_iK_i$ となる有限個の $K_i$ を選べる．
従って
\[
 c(K)\le\sum_i c(K_i)\le\sum_i m(U_i).
\]
$K$ について上限を取ると
\[
 m\!\left(\bigcup_iU_i\right)\le\sum_i m(U_i).
 \tag{4}
\]
これにより $\mu^*$ は零集合・単調性・可算劣加法性を満たす．
また開集合 $U$ 自身を包含する開集合のうち最小のものとして
$U$ を選べるので $\mu^*(U)=m(U)$ である．

次に可測性を確かめる．
互いに交わらないコンパクト集合 $K,L$ は，
互いに交わらない開近傍へ分離できる．
そのための開近傍を一つの開集合 $W$ の内部で取れば，
任意の開集合 $W\supset K\cup L$ に対して
$m(W)\ge c(K)+c(L)$ となる．
従って開集合の内側から近似する方法と
コンパクト集合上の加法性を組み合わせることで，
任意の開集合 $U$ と任意の集合 $A$ に対し
\[
 \mu^*(A)\ge\mu^*(A\cap U)+\mu^*(A\setminus U)
 \tag{5}
\]
が従う．
この操作の詳細を述べると，
$A$ を覆う開集合 $W$ を固定し，
$W$ 内の相対コンパクト開集合による有限被覆を取り，
その閉包を $U$ の内部に入る部分と
$X\setminus U$ の近傍に入る部分に細分する．
閉包が互いに交わらないものについて
$c$ の加法性を適用し，
境界に当たる部分の近傍を縮めて
$m$ の上限と $\mu^*$ の下限を取る．
すると式 (5) を得る．
この分離と正則化の議論は
コンパクト内容からRadon測度を作る
Carath\'eodory拡張手順である．
式 (5) から開集合は $\mu^*$-可測となり，
可測集合全体は $\sigma$-代数なのでBorel集合も可測である．
Borel集合上の $\mu$ は従って可算加法的である．

コンパクト $K$ には相対コンパクト開近傍 $U$ が存在する．
$m(U)\le c(\overline U)<\infty$ だから
$\mu(K)\le\mu(U)<\infty$ である．
外正則性は $\mu^*$ の定義そのものであり，
開集合の内正則性は $\mu(U)=m(U)$ と
$\mu(K)\ge c(K)$ を用いると
\[
 \mu(U)=\sup_{K\subset U}c(K)
 \le\sup_{K\subset U}\mu(K)
 \le\mu(U)
\]
となることから従う．
よって $\mu$ はRadon測度である．
\end{proof}

\begin{theorem}[左Haar測度の存在]\label{thm:haar-exists}
任意の局所コンパクトHausdorff群 $G$ 上に
左Haar測度が存在する．
\end{theorem}
\begin{proof}
定理\ref{thm:content}で不変なコンパクト内容 $c$ を構成し，
補題\ref{lem:content-extension}でRadon測度 $\mu$ へ拡張する．
$c(gK)=c(K)$ なので開集合に対して
\[
 m(gU)=\sup_{K\subset U}c(gK)=m(U)
\]
である．
従って任意の集合 $A$ について
$\mu^*(gA)=\mu^*(A)$ が成立し，
Borel集合上で $\mu(gA)=\mu(A)$ となる．
この測度が零でないことを示す．
$K_0$ は内部が非空で $c(K_0)=1$ なので，
$K_0$ を含む開集合 $U$ では
$\mu(U)=m(U)\ge c(K_0)=1$ である．
よって $\mu$ は非零であり左Haar測度である．
\end{proof}

\section{Haar測度の一意性}

\begin{lemma}[移動不変な可測関数]\label{lem:ergodic}
$\rho$ を全ての非空開集合が正測度となる左Haar測度とし，
$f:G\to\mathbb R$ を有界Borel関数とする．
各 $a\in G$ について
$f(a^{-1}x)=f(x)$ が $\rho$-ほとんど至る所成り立つなら，
$f$ は $\rho$-ほとんど至る所定数である．
\end{lemma}
\begin{proof}
$h\in C_c(G)$ を非負で零でない関数とし，
$H=\int_Gh\,d\rho>0$ と置く．
\[
 F(x)=\int_G h(a)f(a^{-1}x)\,d\rho(a)
\]
を考える．
Fubiniの定理と仮定から
$F(x)=Hf(x)$ が $\rho$-ほとんど至る所成立する．
一方，$a=xy^{-1}$，すなわち
$a=xz$，$z=y^{-1}$ と置いて左不変性を用いると
\[
 F(x)=\int_Gh(xz)f(z^{-1})\,d\rho(z).
\]
$x$ がコンパクト集合内を動くと，
被積分関数の台は共通のコンパクト集合内に入り，
$h(xz)$ は $x$ について一様連続である．
$f$ は有界なので優収束定理によって
$F$ は連続である．
各固定 $b\in G$ について
$F(bx)=Hf(bx)=Hf(x)=F(x)$ が
$\rho$-ほとんど至る所成立する．
左辺と右辺は連続関数であり，
非空開集合は正測度なので，
この等式は全ての $x$ で成立する．
$b,x$ は任意なので $F$ は定数である．
$F=Hf$ がほとんど至る所成立することから
$f$ もほとんど至る所定数となる．
\end{proof}

\begin{theorem}[Haar測度の一意性]\label{thm:haar-unique}
$G$ の二つの左Haar測度 $\mu,\nu$ に対して，
ただ一つの $c>0$ が存在して
\[
 \nu=c\mu
\]
となる．
\end{theorem}
\begin{proof}
$\rho=\mu+\nu$ と置く．
$\rho$ も零でない左不変Radon測度である．
各非空開集合は正測度を持つ：
もし $\rho(U)=0$ なら，$U$ の有限個の左平行移動で
任意のコンパクト集合を覆えるので，
任意のコンパクト集合の測度が零となる．
開集合の内正則性から $\rho=0$ となって矛盾する．

Radon--Nikodymの定理を局所有限な測度に適用し，
\[
 f=\frac{d\mu}{d\rho},\qquad
 g=\frac{d\nu}{d\rho}
\]
を取る．$\mu,\nu\le\rho$ なので
$0\le f,g\le1$ かつ $f+g=1$ が
$\rho$-ほとんど至る所成り立つ．
ここで $\rho$ は必ずしも $\sigma$-有限とは限らないが，
局所コンパクトHausdorff空間のRadon測度に対する
局所的なRadon--Nikodym定理を用いる．
各点の相対コンパクト近傍では全ての測度が有限であり，
重なり上の導関数の一意性で貼り合わせることで
大域的な局所可測導関数が得られる．

左不変性から各 $a\in G$ について
\[
 \int_A f(a^{-1}x)\,d\rho(x)
 =\mu(a^{-1}A)=\mu(A)
 =\int_Af(x)\,d\rho(x)
\]
が全てのBorel集合 $A$ で成り立つ．
従って $f(a^{-1}x)=f(x)$ が
$\rho$-ほとんど至る所成り立つ．
補題\ref{lem:ergodic}から $f$ はほとんど至る所
定数 $\alpha$ である．
同様に $g$ は定数 $\beta$ であり，
$\alpha,\beta\ge0$，$\alpha+\beta=1$ である．
$\mu,\nu$ は零でないので
$\alpha,\beta>0$ となる．
従って
$\mu=\alpha\rho$，$\nu=\beta\rho$ より
$\nu=(\beta/\alpha)\mu$ である．
$\mu$ が零でないので係数は一意である．
\end{proof}

\section{右Haar測度とモジュラー関数}

\begin{definition}
左Haar測度 $\mu$ と $g\in G$ に対し，
$\mu_g(A)=\mu(Ag)$ と置く．
$\mu_g$ も左Haar測度なので，一意性によって
ある $\Delta(g)>0$ が存在して
\[
 \mu(Ag)=\Delta(g)\mu(A)
\]
となる．
写像 $\Delta:G\to(0,\infty)$ を
\emph{モジュラー関数}という．
$\Delta\equiv1$ となる群をunimodular群という．
\end{definition}

\begin{theorem}[モジュラー関数の性質]\label{thm:modular}
$\Delta:G\to(0,\infty)$ は連続群準同型である．
特に
\[
 \Delta(gh)=\Delta(g)\Delta(h),\qquad
 \Delta(g^{-1})=\Delta(g)^{-1}.
\]
またコンパクト群およびAbel群はunimodularである．
\end{theorem}
\begin{proof}
$A$ を $0<\mu(A)<\infty$ となるBorel集合に取る．
定義から
\[
 \mu(Agh)=\Delta(h)\mu(Ag)
 =\Delta(h)\Delta(g)\mu(A)
\]
であり，係数の一意性から積の公式を得る．
逆元の式は $gh=e$ を代入して従う．
連続性を示すため
$f\in C_c(G)$，$f\ge0$，$I=\int f\,d\mu>0$ と取る．
変数変換によって
\[
 \int_G f(xg^{-1})\,d\mu(x)=\Delta(g)\,I
\]
である．
$g$ が一定のコンパクト集合内を動くと，
被積分関数の台は共通のコンパクト集合に入る．
$f(xg^{-1})$ は $g$ について一様連続なので，
優収束定理によって左辺は $g$ に連続に依存する．
$I>0$ だから $\Delta$ も連続である．

$G$ がコンパクトなら $\Delta(G)$ は
$(0,\infty)$ のコンパクト部分群である．
$(0,\infty)$ の非自明な元 $r$ の冪は
一方向に非有界になるので，
そのコンパクト部分群は $\{1\}$ のみである．
従って $\Delta\equiv1$ である．
$G$ がAbelなら右移動と左移動は同じ写像なので，
左不変性から $\mu(Ag)=\mu(gA)=\mu(A)$ となる．
\end{proof}

\begin{proposition}[反転による右Haar測度]\label{prop:right}
$\mu$ を左Haar測度とし
\[
 \nu(A)=\mu(A^{-1})
\]
と置くと，$\nu$ は右Haar測度である．
\end{proposition}
\begin{proof}
反転 $x\mapsto x^{-1}$ は同相写像なので，
$\nu$ はRadon測度である．
任意のBorel集合 $A$ と $g\in G$ に対し
\[
 \nu(Ag)=\mu((Ag)^{-1})
 =\mu(g^{-1}A^{-1})=\mu(A^{-1})=\nu(A)
\]
なので右不変である．
\end{proof}

\section{具体例}

\begin{example}[離散群]
$G$ が離散群なら計数測度は左かつ右Haar測度である．
\end{example}
\begin{proof}
離散空間ではコンパクト集合は有限集合に限られるので
計数測度は有限なコンパクト測度を持ち，
全ての集合が開かつBorelである．
平行移動は集合の全単射なので要素数を保存する．
従って左右に不変なRadon測度となる．
\end{proof}

\begin{example}[Euclid空間]
加法群 $\mathbb R^n$ 上のLebesgue測度はHaar測度であり，
$\Delta\equiv1$ である．
\end{example}
\begin{proof}
Lebesgue測度の平行移動不変性と正則性から
Haar測度の定義を満たす．
群はAbelなのでモジュラー関数は1である．
\end{proof}

\begin{example}[アフィン群]
\[
 G=\{(a,b):a>0,\ b\in\mathbb R\},\qquad
 (a,b)(a',b')=(aa',b+ab')
\]
と置く．
\[
 d\mu(a,b)=a^{-2}\,da\,db
\]
は左Haar測度であり，
\[
 \Delta(a,b)=a^{-1}
\]
である．
\end{example}
\begin{proof}
左移動
$(a,b)\mapsto(a_0a,b_0+a_0b)$ の
Jacobi行列式は $a_0^2$ である．
従って
$(a_0a)^{-2}a_0^2\,da\,db
 =a^{-2}\,da\,db$
となり，$\mu$ は左不変である．
右移動
$(a,b)\mapsto(aa_0,b+ab_0)$ の
Jacobi行列式は $a_0$ であるから
\[
 (aa_0)^{-2}a_0\,da\,db
 =a_0^{-1}a^{-2}\,da\,db.
\]
よって $\mu(A(a_0,b_0))=a_0^{-1}\mu(A)$ となり，
モジュラー関数の表示を得る．
\end{proof}

\begin{remark}
局所コンパクト性は測度の存在に重要である．
特にコンパクト群のHaar測度は全測度1へ正規化できるが，
非コンパクト群上には有限で非零な左不変Radon測度は存在しない．
実際全測度が有限なら，
正測度の相対コンパクト開集合を互いに交わらないよう
無限回平行移動できる非コンパクト群では矛盾する．
\end{remark}

\section{到達課題の完全解答}
\noindent\textbf{円周群のHaar積分．}
\begin{proposition}
円周群 $S^1$ 上の規格化Haar積分は，
連続関数 $f$ に対して
\[
\int_{S^1}f(z)\,d\mu(z)
=\frac1{2\pi}\int_0^{2\pi}f(e^{it})\,dt
\]
で与えられる．
\end{proposition}
\begin{proof}
右辺は正の線形汎関数で，
$f=1$ の積分は1である．
群の任意の元 $e^{i\theta}$ による左移動に対して，
周期性と変数変換 $s=t+\theta$ から積分値は不変となる．
正規化したHaar測度の一意性により
この積分表示がHaar積分を与える．
\end{proof}
\begin{corollary}
任意の整数 $n\ne0$ に対し
\[
\int_{S^1}z^n\,d\mu(z)=0
\]
が成立する．
\end{corollary}
\begin{proof}
$(2\pi)^{-1}\int_0^{2\pi}e^{int}\,dt=0$ である．
\end{proof}
\noindent\textbf{到達確認．}
$\int_{S^1}|1+z|^2\,d\mu(z)$ を計算せよ．
\par\noindent\textbf{解答．}
$|1+z|^2=(1+z)(1+\bar z)=2+z+\bar z$．
非自明な指標の積分は零なので値は2である．
\end{document}